\documentclass[fontset=none,zihao=-4]{Notes}

\makeatletter
\DeclareRobustCommand{\em}{%
  \@nomath\em \if b\expandafter\@car\f@series\@nil
  \normalfont \else \bfseries \fi}
\makeatother

\usepackage{tikz-cd,wrapstuff}
\usepackage{fixdif,siunitx,tikz,nicematrix}

\usetikzlibrary{hobby,calc,arrows}
\usetikzlibrary{positioning}
\usetikzlibrary{decorations.markings}
\usetikzlibrary{decorations.pathreplacing}

\input{font.def}

\DeclareMathOperator\Int{Int}
\DeclareMathOperator\supp{supp}
\DeclareMathOperator\im{im}
\DeclareMathOperator\End{End}
\DeclareMathOperator\Ann{Ann}
\DeclareMathOperator\Hom{Hom}
\DeclareMathOperator\diag{diag}
\DeclareMathOperator\Or{O}
\DeclareMathOperator\rank{rank}
\DeclareMathOperator\Ob{Ob}
\DeclareMathOperator\Ad{Ad}
\DeclareMathOperator\ad{ad}
\DeclareMathOperator\id{id}
\DeclareMathOperator\SO{SO}
\DeclareMathOperator\SU{SU}
\DeclareMathOperator\GL{GL}
\DeclareMathOperator\SL{SL}

\newcommand{\cat}[1]{\mathsf{#1}}

\newcommand{\inn}[1]{\left\langle#1\right\rangle}
\newcommand{\lie}[1]{\mathfrak{#1}}
\newcommand{\norm}[1]{\left\lVert#1\right\rVert}
\newcommand{\spa}[1]{\operatorname{span}\left(#1\right)}

\makeatletter
\DeclareRobustCommand\bigop[2][1]{%
  \mathop{\vphantom{\sum}\mathpalette\bigop@{{#1}{#2}}}\slimits@
}
\newcommand{\bigop@}[2]{\bigop@@#1#2}
\newcommand{\bigop@@}[3]{%
  \vcenter{%
    \sbox\z@{$#1\sum$}%
    \hbox{\resizebox{\ifx#1\displaystyle#2\fi\dimexpr\ht\z@+\dp\z@}{!}{$\m@th#3$}}%
  }%
}
\makeatother

\newcommand{\fprod}{\DOTSB\bigop[1.05]{\ast}}

\tikzcdset{
  arrow style=tikz,
  diagrams={>={Straight Barb[scale=0.8]}}
}

\tikzset{
  point/.style={
    circle,fill,inner sep=0pt,minimum width=5pt
  }
}

\usepackage[subscriptcorrection,nofontinfo,mtpbb,mtpfrak,mtpscr]{mtpro2}

\newcommand{\mat}[1]{\mathbold{#1}}

\usepackage{biblatex}
% \addbibresource{ref.bib}
% \begin{filecontents}{ref.bib}
% @book{kirillov2008introduction,
%   title={An introduction to Lie groups and Lie algebras},
%   author={Kirillov, Alexander A},
%   volume={113},
%   year={2008},
%   publisher={Cambridge University Press}
% }

% @article{notes,
%   author = {Gallier, Jean},
%   year = {2009},
%   month = {02},
%   pages = {},
%   title = {Notes on Spherical Harmonics and Linear Representations of Lie Groups}
% }
% \end{filecontents}

\setlist[enumerate]{nosep,label=(\arabic*)}
\setlist[itemize]{nosep}

\title{\sffamily 强化学习的数学原理}
\author{Eliauk}



\begin{document}

\maketitle

\tableofcontents

\section{Markov 决策过程}

\subsection{带折扣的 Markov 决策过程}

在强化学习中，agent 与环境的交互被建模为一个带折扣的 Markov 决策过程
(Markov Decision Process, MDP)：
\[
  M=(\mathcal S, \mathcal A,  P, r, \gamma,\mu).
\]
其中各项的定义如下：
\begin{itemize}
  \item 状态空间 $\mathcal S$，可能有限或者无限。一般来说，要求 $\mathcal S$ 是一个可测空间，
  配备一个 $\sigma$-代数 $\mathscr{S}$。为了数学上的便利，我们要求 $\mathcal S$ 是
  有限或者可数无限的。
  \item 动作空间 $\mathcal A$，同样可能有限或者无限。一般来说，要求 $\mathcal A$ 是一个可测空间，
  配备一个 $\sigma$-代数 $\mathscr{A}$。为了数学上的便利，我们要求 $\mathcal A$ 是
  有限的。
  \item 转移函数 $P:(\mathcal S\times \mathcal A)\times \mathscr{S}\to [0,1]$，满足：对于任意 $(s,a)\in \mathcal S\times \mathcal A$，$P(\cdot | s,a)$ 是 $\mathcal S$ 上的一个概率测度。
  我们也记 $P_{s,a} := P(\cdot | s,a)$。
  \item 奖励函数 $r: \mathcal S \times \mathcal A \to [0,1]$，表示在状态 $s$ 下执行动作 $a$ 所获得的即时奖励。
  $r$ 视为 $\mathcal S \times \mathcal A$ 上的一个随机变量，且要求 $r$ 是有界的。
  \item 折扣因子 $\gamma\in [0,1)$，定义了问题的有效视野。
  \item 初值分布 $\mu$，是 $\mathcal S$ 上的一个概率测度，表示系统的初始状态分布。
在许多情况下，我们假设初值状态是固定的 $s_0$，即 $\mu = \delta_{s_0}$。
\end{itemize}

\subsubsection{目标、策略和价值}

\paragraph{策略} 给定 MDP $M=(\mathcal S, \mathcal A,  P, r, \gamma,\mu)$，agent 遵从下面的原则和环境进行交互：
agent 从某个状态 $s_0\sim \mu$ 开始；在每个时间 $t=0,1,2,\dots$ 处，agent 首先观测当前状态 $s_t$，
然后选择一个动作 $a_t\in \mathcal A$，获得一个即时奖励 $r_t=r(s_t,a_t)$，最后再观测下一个状态 
$s_{t+1}\sim P(\cdot|s_t,a_t)$。

在一个决策时间 $t$ 时，约定历史信息为截至 $t$ 的可观测的状态和动作的序列，即
\[
  \tau_t:=(s_0,a_0,s_1,a_1,\dots,s_{t-1},a_{t-1},s_t).
\]
注意到 $\tau_t$ 中不包含 $a_t$，因为 $a_t$ 是在历史信息 $\tau_t$ 的基础上选择的。

在最一般的情况下，一个策略指的是 agent 自适应地根据历史信息来选择动作的规则。准确地说，
一个策略是一个映射 $\pi:\mathcal H\times \mathscr A \to [0,1]$，其中 $\mathcal H$ 包含所有有限长的历史信息。
在时间 $t$ 处，动作采样于 $ a_t\sim \pi(\cdot|\tau_t)$。

一个稳定策略 $\pi$ 是一个映射 $\pi:\mathcal S\times \mathscr A \to [0,1]$，满足：对于任意 $s\in \mathcal S$，$\pi(\cdot|s)$ 是 $\mathcal A$ 上的一个概率测度。
也即，agent 选取的动作采样于 $a_t\sim \pi(\cdot|s_t)$，只依赖于当前状态 $s_t$。
一个确定性的稳定策略指的是映射 $\pi:\mathcal S \to \mathcal A$，满足 $\pi(s)\in \mathcal A$，即在状态 $s$ 下总是选择动作 $\pi(s)$。

\paragraph{价值} 我们现在定义一般策略的价值。对于一个策略 $\pi$ 和初始状态 $S_0=s$，我们定义价值函数 $V^\pi_M: \mathcal S \to \mathbb{R}$ 为
\[
  V_M^\pi(s)= \mathbb{E}\biggl[
    \sum_{t=0}^\infty \gamma^t r(S_t,A_t) \biggm| S_0=s
  \biggr].
\]
这里的期望相对于轨迹的随机性(依赖于状态转移和策略 $\pi$)。因为 $r(s,a)\in [0,1]$，所以
$ 0<V_M^\pi(s)\leq \sum_{t=0}^\infty \gamma^t = \frac{1}{1-\gamma}$。

类似地，定义动作价值函数 $Q^\pi_M:\mathcal S\times \mathcal A \to \mathbb{R}$ 为
\[
  Q_M^\pi(s,a) = \mathbb{E}\biggl[
    \sum_{t=0}^\infty \gamma^t r(S_t,A_t) \biggm| S_0=s, A_0=a
  \biggr].
\]
显然 $Q_M^\pi(s,a)$ 也有上界 $\frac{1}{1-\gamma}$。

\paragraph{目标} 给定状态 $s$，agent 的目标是寻找策略 $\pi$ 使得最大化价值：
\begin{equation}
  \sup_\pi V_M^\pi(s),
\end{equation}
其中上确界针对所有可能的(非稳定的或者随机的)策略。我们将看到，存在一个稳定的确定性策略对于所有初始状态 $s$ 
都是最优的。

为了简化记号，我们有时候会省略 MDP 的下标 $M$，直接写成 $V^\pi$ 和 $Q^\pi$。

\begin{example}[导航]
  导航也许是最简单的强化学习的例子。状态是 agent 的当前位置。动作为东南西北四个移动方向。
  此时状态转移函数是确定性的。
  agent 有一个目标状态 $g$，当 agent 到达目标状态 $g$ 时获得奖励 $1$，其他时候奖励为 $0$。
  因为 $\gamma <1$，所以存在一个激励使得 agent 尽快地到达目标状态 $g$。
  因此，最优的行为是找到从起点到目标状态 $g$ 的最短路径。如果策略在 $d$ 步到达目标并且保
  持在目标状态，那么价值函数的值为 $\gamma^d$。
\end{example}

\begin{example}[策略游戏]
  这是一类很受欢迎的应用场景，强化学习已经在围棋、国际象棋以及各种形式的扑克中达到了人类水平的表现。通常，状态
  是当前的游戏局面，动作是合法的走法，而奖励则是最终的胜负结果（或者是更详细的得分）。从技术上讲，这些属于
  多 agent 环境，但许多成功的方法使用的技术并不会显式地对多个 agent 建模。
\end{example}

\subsubsection{稳定策略的 Bellman 一致性方程}

稳定策略满足下面的一致性条件：

\begin{lemma}
  假设 $\pi$ 是稳定策略，那么 $V^\pi$ 和 $Q^\pi$ 满足下面的 Bellman 一致性方程：
  对于任意 $s\in \mathcal S$ 和 $a\in \mathcal A$，有
  \begin{align*}
    V^\pi(s)&=\mathbb{E}_{a\sim \pi(\cdot|s)} [Q^\pi(s,a)],\\
    Q^\pi(s,a) &= r(s,a) + \gamma \mathbb{E}_{s'\sim P(\cdot|s,a)} [V^\pi(s')].
  \end{align*}
  特别地，如果 $\pi$ 是确定性的，那么 $V^\pi(s)=Q^\pi(s,\pi(s))$。
\end{lemma}
\begin{proof}
  直接利用条件期望的塔式法则，我们有
  \begin{equation*}
    \mathbb{E}\biggl[
      \mathbb{E}\Bigl[
        \sum_{t=0}^\infty \gamma^t r(S_t,A_t)\Bigm| S_0,A_0
      \Bigr]\biggm| S_0
    \biggr]= \mathbb{E}\biggl[
      \sum_{t=0}^\infty \gamma^t r(S_t,A_t) \biggm| S_0
      \biggr],
  \end{equation*}
  右端是 $V^\pi$，左端是 $\mathbb{E}_{a\sim \pi(\cdot|S_0)}[Q^\pi(S_0,a)]$。
  对于第二点，有
  \begin{align*}
    Q^\pi&=\mathbb{E}\biggl[
      r(S_0,A_0) + \sum_{t=1}^\infty \gamma^t r(S_t,A_t) \biggm| S_0,A_0
    \biggr]\\
    &=r(S_0,A_0)+\gamma \mathbb{E}\biggl[
      \sum_{t=0}^\infty \gamma^t r(S_{t+1},A_{t+1}) \biggm| S_0,A_0
    \biggr]\\
    &= r(S_0,A_0) + \gamma \mathbb{E}\biggl[
      \mathbb{E}\Bigl[
        \sum_{t=0}^\infty \gamma^t r(S_{t+1},A_{t+1}) \Bigm| S_1
      \Bigr] \biggm| S_0,A_0
    \biggr]\\
    &= r(S_0,A_0) + \gamma \mathbb{E}_{s'\sim P(\cdot|S_0,A_0)}[V^\pi(s')].\qedhere
  \end{align*}
\end{proof}

换一个视角，把 $V^\pi$ 视为长度 $|\mathcal S|$ 的向量，$Q^\pi$ 视为长度 $|\mathcal S|\cdot |\mathcal A|$ 的向量。
记符号 $P$ 为转移函数 $P$ 的矩阵表示，即一个 $|\mathcal S|\cdot |\mathcal A|$ 行 $|\mathcal S|$ 列的矩阵，
满足 $P_{(s,a),s'}=P(s'|s,a)$。
同时记 $P^\pi$ 是稳定策略 $\pi$ 诱导的状态-动作对的转移矩阵：
\[
  P^\pi_{(s,a),(s',a')} = P(s'|s,a)\pi(a'|s').
\]
特别的，对于确定性策略，有
\[
  P^\pi_{(s,a),(s',a')} = \begin{cases}
    P(s'|s,a) & \text{if $a'=\pi(s')$},\\
    0 & \text{otherwise}.
  \end{cases}
\]
在这个记号下，可以直接验证得到：
\[
  Q^\pi=r+\gamma PV^\pi=r+\gamma P^\pi Q^\pi.
\]

\begin{corollary}
  假设 $\pi$ 是稳定策略，那么 
  \begin{equation}
    Q^\pi=(I-\gamma P^\pi)^{-1}r,
  \end{equation}
  其中 $I$ 是 $|\mathcal S|\cdot |\mathcal A|$ 维的单位矩阵。
\end{corollary}
\begin{proof}
  只需要说明 $I-\gamma P^\pi$ 是可逆的。注意到对于任意非零的向量 $x\in \mathbb{R}^{|\mathcal S| |\mathcal A|}$，有
  \begin{equation*}
    \| (I-\gamma P^\pi)x\|_\infty \geq \|x\|_\infty - \gamma \|P^\pi x\|_\infty\geq 
    \|x\|_\infty - \gamma \|x\|_\infty = (1-\gamma)\|x\|_\infty >0.
  \end{equation*}
  其中第一步利用了三角不等式，第二步利用了 $P^\pi$ 是一个转移矩阵，所以 $\|P^\pi x\|_\infty \leq \|x\|_\infty$。
\end{proof}

\subsubsection{Bellman 最优性方程}

带折扣的 MDP 的一个非常重要和方便的性质是，存在一个稳定和确定性的策略使得对于所有的 $s \in \mathcal S$
都能最大化 $V^\pi(s)$。

\begin{theorem}
  令 $\Pi$ 是所有策略(可能是非稳定的或者随机的)的集合。定义
  \[
    V^*(s) = \sup_{\pi\in \Pi} V^\pi(s),\quad Q^*(s,a) = \sup_{\pi\in \Pi} Q^\pi(s,a).
  \]
  那么存在一个稳定的确定性策略 $\pi$，使得对于所有 $s\in \mathcal S$ 和 $a\in \mathcal A$，满足 
  \[
    V^\pi(s) = V^*(s),\quad Q^\pi(s,a)=Q^*(s,a).
  \]
  我们把这样的策略称为最优策略。
\end{theorem}
\begin{proof}
  我们首先证明：在考虑一步历史 $(S_0,A_0,S_1)=(s,a,s')$ 的条件下，最大的未来折扣回报
  不依赖于 $(s,a)$ 的值。准确地说，我们证明下面的式子：
  \begin{equation}\label{eq:optimality}
    \sup_{\pi\in \Pi} \mathbb{E}\biggl[
      \sum_{t=1}^\infty \gamma^t r(S_t,A_t) \biggm| S_0=s, A_0=a, S_1=s'
    \biggr]=\gamma V^*(s').
  \end{equation}
  固定任意策略 $\pi$。定义一个“偏移”策略 $\pi_{(s,a)}$ 如下：对于 $t\geq 0$ 和历史 $(s_0,a_0,\dots,s_t)$，有
  \begin{align*}
    &\pi_{(s,a)}(A_t=b\mid S_0=s_0,A_0=a_0,\dots,S_t=s_t)\\&=
    \pi(A_{t+1}=b\mid S_0=s,A_0=a,S_1=s_0,A_1=a_0,\dots,S_{t+1}=s_t).
  \end{align*}
  那么在已知 $(S_0,A_0,S_1)=(s,a,s')$ 的条件下策略 $\pi$ 的表现，与
  已知 $S_0=s'$ 的条件下策略 $\pi_{(s,a)}$ 的表现完全一样。也就是说，有
  \[
    \mathbb{E}^\pi \biggl[
      \sum_{t=1}^\infty \gamma^t r(S_t,A_t) \biggm| (S_0,A_0,S_1)=(s,a,s')
    \biggr] = \gamma \mathbb{E}^{\pi_{(s,a)}} \biggl[
      \sum_{t=0}^\infty \gamma^t r(S_t,A_t) \biggm| S_0=s'
    \biggr]=\gamma V^{\pi_{(s,a)}}(s').
  \]
  此外，对于每个 $(s,a)$，集合 $\{\pi_{(s,a)} : \pi \in \Pi\}$ 与 $\Pi$ 是相同的。
  因此，有
  \[
    \sup_{\pi\in \Pi} \mathbb{E}^\pi \biggl[
      \sum_{t=1}^\infty \gamma^t r(S_t,A_t) \biggm| (S_0,A_0,S_1)=(s,a,s')
    \biggr] = \gamma \sup_{\pi\in \Pi} V^{\pi_{(s,a)}}(s') = \gamma V^*(s').
  \]
  这就证明了 \eqref{eq:optimality} 式。


\end{proof}









% \section{系统动力学的公理化构造}

% 本节我们构建强化学习的系统框架，即马尔可夫决策过程（Markov Decision Process, MDP）。
% 我们将从最基本的概念出发，逐步建立起 MDP 的数学定义。

% MDP 指的是一个六元组 $(S, A, \mathcal P, r,\rho_0, \gamma)$，其中：
% \begin{enumerate}
%   \item \emph{状态空间}。系统所有可能的状态的集合，记为 $S$。默认
%   是一个可测空间，配备一个 $\sigma$-代数 $\mathcal{S}$。
%   \item \emph{动作空间}。系统所有可能的动作的集合，记为 $A$。同样是一个可测空间，配备一个 $\sigma$-代数 $\mathcal{A}$。
%   为了描述系统的交互，还会考虑乘积空间 $S \times A$，配备 $\sigma$-代数 $\mathcal{S} \otimes \mathcal{A}$。
%   \item \emph{状态转移核}。状态转移核 $\mathcal P$ 指的是一个从 $S\times A$ 到 
%   $S$ 的转移概率，也即映射
%   \[
%     \mathcal P: (S\times A)\times \mathcal{S} \to [0,1],
%   \]
%   满足：对于任意 $(s,a)\in S\times A$，$\mathcal P(s,a,\cdot)$ 是 $S$ 上的一个概率测度；
%   对于任意 $E\in \mathcal{S}$，$\mathcal P(\cdot,\cdot,E)$ 是 $S\times A$ 上的一个可测函数。
%   通常，我们记 $\mathcal P(E|s,a) := \mathcal P(s,a,E)$，表示在状态 $s$ 下执行动作 $a$ 后状态集 $E$ 
%   出现的概率。
%   \item \emph{奖励函数}。奖励函数 $r$ 是 $S\times A$ 上的有界实值随机变量。
%   \item \emph{初始状态分布}。$\rho_0$ 是 $S$ 上的一个概率测度，表示系统的初始状态分布。
%   \item \emph{折扣因子}。$\gamma\in [0,1)$ 是一个实数，表示未来奖励的折扣率。
% \end{enumerate}

% 在此基础上，我们需要优化的目标是一个策略 $\pi_\theta$，它是从 $S$ 到 $A$ 的转移概率，所以
% 满足：对于任意 $s\in S$，$\pi_\theta(s,\cdot)$ 是 $A$ 上的一个概率测度；对于任意 $E\in \mathcal{A}$，$\pi_\theta(\cdot,E)$ 是 $S$ 上的一个可测函数。
% 同样记 $\pi_\theta(E|s) := \pi_\theta(s,E)$，表示在状态 $s$ 下选择动作集 $E$ 的概率。
% $\theta$ 表示策略的参数。

% \section{无穷维轨迹空间与目标泛函}

% 在强化学习中，系统的完整状态是一条无限长的序列，我们需要构造一个容纳
% 所有可能轨迹的空间。

% 对于每个时刻 $t\in \mathbb{N}$，系统的状态是一个二元组 $(s_t,a_t)\in S \times A$，表示在时刻 $t$ 的状态和动作。
% 因此，系统的完整状态可以表示为一个无限长的序列 $\tau = ((s_0,a_0),(s_1,a_1),\ldots)$，其中每个 $(s_t,a_t)\in S \times A$。
% 因此我们的状态空间应该是
% \[
%   \Omega=(S\times A)^{\mathbb{N}}=\prod_{t=0}^\infty S\times A,
% \]
% 其配备自然的 $\sigma$-代数 $\mathcal{F}$，由所有有限维的柱集生成。

% 接下来我们要给 $(\Omega,\mathcal F)$ 赋予一个概率测度 $\mathbb P_\theta$。
% 现在我们拥有的是：
% \begin{itemize}
%   \item 初始概率测度 $\rho_0(\d s_0)$。
%   \item 策略核 $\pi_\theta(\d a_t|s_t)$，即给定状态输出动作的概率测度。
%   \item 系统转移核 $\mathcal P(\d s_{t+1}|s_t,a_t)$，即
%   给定当前时刻状态和动作输出下一个时刻状态的概率测度。
% \end{itemize}
% 我们定义一个从 $S\times A$ 到 $S\times A$ 的联合转移核
% \[
%   K_\theta(\d s_{t+1},\d a_{t+1}|s_t,a_t) \coloneq \mathcal P(\d s_{t+1}|s_t,a_t)\pi_\theta(\d a_{t+1}|s_{t+1}).
% \]
% 换句话说，给定 $s_t,a_t$，$K_\theta(\cdot| s_t,a_t)$ 是积测度 $\mathcal P(\cdot|s_t,a_t)\otimes \pi_\theta(\cdot|s_{t+1})$，
% 根据 Ionescu-Tulcea 定理，空间 $(\Omega,\mathcal F)$ 上存在唯一的
% 概率测度 $\mathbb P_\theta$，使得对于任意时间 $T\in \mathbb{N}$
% 和可测集 $B_0,B_1,\dots,B_T\in \mathcal S$ 以及 $C_0,C_1,\dots,C_T\in \mathcal A$，满足
% \begin{align*}
%   &\mathbb P_\theta\left(
%     (B_0\times C_0)\times (B_1\times C_1)\times \cdots \times (B_T\times C_T) \times (S\times A)^{\mathbb{N}}
%   \right) \\
%   &= \int_{B_0}\rho_0(\d s_0)
%   \int_{C_0}\pi_\theta(\d a_0|s_0)
%   \int_{B_1\times C_1} K_\theta(\d s_1,\d a_1|s_0,a_0)
%   \cdots \int_{B_T\times C_T} K_\theta(\d s_T,\d a_T|s_{T-1},a_{T-1})\\
%   &=\int_{B_0\times C_0}\int_{B_1\times C_1} \cdots \int_{B_T\times C_T} \rho_0(\d s_0)
%   \pi_\theta(\d a_0|s_0)
%   \prod_{t=0}^{T-1} K_\theta(\d s_{t+1},\d a_{t+1}|s_{t},a_{t})
% \end{align*}
% 这也表明 $\mathbb{P}_\theta$ 有概率密度函数
% \begin{equation}\label{eq:prob-density}
%   p_\theta(\tau) = \prod_{t=-1}^\infty K_\theta(s_{t+1},a_{t+1}|s_{t},a_{t}).
% \end{equation}
% 这里为了方便，记 $K_{\theta}(\d s_0,\d a_0|s_{-1},a_{-1}) = \rho_0(\d s_0)\pi_\theta(\d a_0|s_0)$。

% 有了概率空间 $(\Omega,\mathcal F,\mathbb P_\theta)$，我们就可以定义强化学习的目标函数了。
% 定义投影函数 $X_t:\Omega \to S \times A$ 为 $X_t(\tau) = (s_t,a_t)$，其中 $\tau = ((s_0,a_0),(s_1,a_1),\ldots) \in \Omega$。
% 我们定义单步奖励是复合函数 $r_t=r\circ X_t:\Omega \to \mathbb{R}$，即
% \[
%   r_t(\tau)=r(s_t,a_t).
% \]
% $r_t$ 是 $\Omega$ 上的有界实值随机变量。
% 通过折扣因子 $\gamma$，我们定义整条轨迹的累计奖励为无穷级数
% \begin{equation}\label{eq:reward}
%   R(\tau)=\sum_{t=0}^\infty \gamma^t r_t(\tau).
% \end{equation}
% 由于 $r_t$ 的有界性，$R(\tau)$ 是绝对收敛的，并且关于 $\tau$ 是一致收敛的，
% 所以 $R$ 也是 $\Omega$ 上的一个有界实值随机变量。
% 也就是说，$R(\tau)\in L^1(\Omega,\mathcal F,\mathbb P_\theta)$。
% 最后，强化学习的\emph{目标泛函} $J:\Theta\to \mathbb{R}$ 定义为
% 随机变量 $R$ 的期望值，即
% \begin{equation}\label{eq:objective-functional}
%   J(\theta)= \mathbb{E}[R]=\int_\Omega R(\tau)\mathbb P_\theta(\d \tau).
% \end{equation}
% 可以看到，强化学习的奖励并不依赖于参数 $\theta$。
% 参数 $\theta$ 唯一控制的变量是概率测度 $\mathbb P_\theta$。
% 因此强化学习优化的本质是在参数空间 $\Theta$ 上不断改变测度 $ \mathbb{P}_\theta$，
% 把更多的“概率质量”集中到 $R$ 较大的轨迹集合上。

% \section{策略梯度}

% 为了书写方便，我们先对机器学习文献中常见的期望写法进行一些符号上的说明。
% 对于随机变量 $f:\Omega \to \mathbb{R}$，记号 $\mathbb{E}_{x\sim \mathbb{P}}[f(x)]$
% 表示 $f$ 关于测度 $\mathbb{P}$ 的期望值，即
% \[
%   \mathbb{E}_{x\sim \mathbb{P}}[f(x)] = \int_\Omega f(x)\mathbb{P}(\d x).
% \] 
% 有时候也简写成 $\mathbb{E}_x[f(x)]$。对于带有转移核的情况，
% 设两个可测空间 $(E,\mathcal{E})$ 和 $(F,\mathcal{F})$，以及一个从 $E$ 到 $F$ 的转移核 $K:E\times \mathcal{F}\to [0,1]$，
% $\mathbb{P}$ 是 $E$ 上的概率测度，$f$ 是 $E\times F$ 上的随机变量，那么
% 记号 $\mathbb{E}_{x\sim \mathbb{P},y\sim K(\cdot|x)}[f(x,y)]$ 表示
% \[
%   \mathbb{E}_{x\sim \mathbb{P},y\sim K(\cdot|x)}[f(x,y)] = \int_E \int_F f(x,y) K(\d y|x)\mathbb{P}(\d x).   
% \]
% 有时候也简写成 $\mathbb{E}_{x\sim \mathbb{P},y\sim K}[f(x,y)]$ 或者 $\mathbb{E}_{x,y}[f(x,y)]$。

% 回顾 \eqref{eq:objective-functional} 式，结合累计奖励 \eqref{eq:reward} 和
% 概率密度 \eqref{eq:prob-density}，我们可以把目标函数 $J(\theta)$ 展开成如下形式：
% \[
%   J(\theta)=\int_\Omega \left(
%     \sum _{k=0}^\infty \gamma^k r(s_k,a_k)
%   \right) p_\theta(\tau) \d \tau,
% \]
% 对 $\theta$ 求梯度，交换积分和梯度算子，利用恒等变形 $\nabla_\theta p_\theta(\tau) = p_\theta(\tau)\nabla_\theta \log p_\theta(\tau)$，我们得到
% \[
%   \nabla_\theta J(\theta) = \int_\Omega \left(
%     \sum _{k=0}^\infty \gamma^k r(s_k,a_k)
%   \right) p_\theta(\tau)\nabla_\theta \log p_\theta(\tau) \d \tau
%   =\mathbb{E}_{\tau\sim \mathbb{P}_\theta}\left[
%     \nabla_\theta \log p_\theta(\tau)
%     \sum _{k=0}^\infty \gamma^k r(s_k,a_k)
%   \right].
% \]
% 再回顾概率密度 \eqref{eq:prob-density}，
% 表达式中只有 $\pi_\theta(a_t|s_t)$ 依赖于 $\theta$，所以
% \[
%   \nabla_\theta \log p_\theta(\tau) = \sum_{t=0}^\infty \nabla_\theta \log \pi_\theta(a_t|s_t),
% \]
% 代入到上式中，我们就得到了强化学习的\emph{原始策略梯度表达式}：
% \begin{equation}\label{eq:policy-gradient}
%   \nabla_\theta J(\theta) = \mathbb{E}_{\tau\sim \mathbb{P}_\theta}\left[
%     \sum_{t=0}^\infty\sum_{k=0}^\infty \gamma^k r(s_k,a_k) \nabla_\theta \log \pi_\theta(a_t|s_t)
%   \right].  
% \end{equation}

% 原始策略梯度的表达式极难处理，包含了二重无穷级数，
% 我们需要利用时间上的因果性来进行一些简化。

% 记 $S_k:\Omega\to S$ 和 $A_k:\Omega\to A$ 分别为时刻 $k$ 的状态和动作的投影函数，即 $S_k(\tau)=s_k$ 和 $A_k(\tau)=a_k$。
% 对于每个时刻 $t$，我们考虑历史轨迹生成子 $\sigma$-代数 $\mathcal H_t$：
% \[
%   \mathcal H_t = \sigma\left(
%     S_0,A_0,S_1,A_1,\dots,S_{t-1},A_{t-1},S_t
%   \right).
% \]
% 利用条件期望的塔式法则，我们有
% \[
%   \mathbb{E}_\tau\left[
%     \nabla_\theta\log \pi_\theta(a_t|s_t)\sum_{k=0}^{t-1}\gamma^k r_k
%   \right]= \mathbb{E}\left[
%     \mathbb{E}\left[
%       \nabla_\theta\log \pi_\theta(A_t|S_t)\sum_{k=0}^{t-1}\gamma^k r_k
%       \middle| \mathcal H_t
%     \right]
%   \right],
% \]
% 内层期望中，$\sum_{k=0}^{t-1}\gamma^k r_k$ 是 $\mathcal H_t$-可测的，
% 根据条件期望的性质，我们可以把它往外提一层，得到
% \[
%   \mathbb{E}\left[
%     \nabla_\theta\log \pi_\theta(A_t|S_t)\sum_{k=0}^{t-1}\gamma^k r_k
%     \middle| \mathcal H_t
%   \right] = \left(\sum_{k=0}^{t-1}\gamma^k r_k\right) \mathbb{E}\left[
%     \nabla_\theta\log \pi_\theta(A_t|S_t)
%     \middle| \mathcal H_t
%   \right].
% \]
% 其中 $\pi_\theta(A_t|S_t)$ 只有 $A_t$ 是未知信息，所以有
% \begin{equation}\label{eq:zero-mean}
%   \begin{aligned}
%     \mathbb{E}[
%       \nabla_\theta\log \pi_\theta(A_t|S_t)
%       | \mathcal H_t]&=\mathbb{E}[
%       \nabla_\theta\log \pi_\theta(A_t|S_t)|S_t]
%       =\int_A 
%       \nabla_\theta\log \pi_\theta(a_t|S_t)\pi_\theta(\d a_t|S_t)\\
%       &=\int_A \frac{\nabla_\theta \pi_\theta(a_t|S_t)}{\pi_\theta(a_t|S_t)}\pi_\theta( a_t|S_t)
%       \d a_t=\int_A \nabla_\theta \pi_\theta(a_t|S_t) \d a_t\\
%       &=\nabla_\theta \int_A \pi_\theta(a_t|S_t) \d a_t=\nabla_\theta 1=0.
%   \end{aligned}
% \end{equation}
% 于是我们有
% \[
%   \mathbb{E}_\tau\left[
%     \nabla_\theta\log \pi_\theta(a_t|s_t)\sum_{k=0}^{t-1}\gamma^k r_k
%   \right]= 0.
% \]
% 所以 \eqref{eq:policy-gradient} 中关于 $k$ 的求和可以截断到 $t$，即
% \begin{equation}\label{eq:policy-gradient-simplified}
%   \nabla_\theta J(\theta) = \mathbb{E}_{\tau\sim \mathbb{P}_\theta}\left[
%     \sum_{t=0}^\infty\sum_{k=t}^\infty \gamma^k r(s_k,a_k) \nabla_\theta \log \pi_\theta(a_t|s_t)
%   \right].
% \end{equation}

% 把 $t$ 的求和与期望交换并提取因子 $\gamma^t$，得到
% \[
%   \nabla_\theta J(\theta) = \sum_{t=0}^\infty \gamma^t\mathbb{E}_{\tau\sim \mathbb{P}_\theta}\left[
%     \nabla_\theta \log \pi_\theta(a_t|s_t) \sum_{k=t}^\infty \gamma^{k-t} r(s_k,a_k)
%   \right].
% \]
% 继续利用条件期望的塔式法则，我们有
% \begin{align*}
%   &\mathbb{E}_{\tau\sim \mathbb{P}_\theta}\left[
%     \nabla_\theta \log \pi_\theta(a_t|s_t) \sum_{k=t}^\infty \gamma^{k-t} r(s_k,a_k)
%   \right]\\
%   &= \mathbb{E}\left[
%     \mathbb{E}\left[
%       \nabla_\theta \log \pi_\theta(A_t|S_t) \sum_{k=t}^\infty \gamma^{k-t} r_k
%       \middle| S_t,A_t
%     \right]
%   \right]\\
%   &=\mathbb{E}\left[
%     \nabla_\theta \log \pi_\theta(A_t|S_t) \mathbb{E}\left[
%       \sum_{k=t}^\infty \gamma^{k-t} r_k
%       \middle| S_t,A_t
%     \right]
%   \right],
% \end{align*}
% 第二个等号是因为 $\nabla_\theta \log \pi_\theta(A_t|S_t)$ 是 $\sigma(S_t,A_t)$-可测的。
% 记\emph{动作价值函数}为
% \[
%   Q^{\pi_\theta} = \mathbb{E}\left[
%     \sum_{k=t}^\infty \gamma^{k-t} r_k
%     \middle| S_t,A_t
%   \right],
% \]
% 其表示在状态 $S_t$ 下执行动作 $A_t$ 后，未来累计奖励的条件期望，是 $(S_t,A_t)$ 的函数。
% 于是策略梯度表达式变为
% \[
%   \nabla_\theta J(\theta) = \sum_{t=0}^\infty \gamma^t\mathbb{E}_{\tau\sim \mathbb{P}_\theta}\left[
%     \nabla_\theta \log \pi_\theta(a_t|s_t) Q^{\pi_\theta}(s_t,a_t)
%   \right].
% \]
% % 由于期望内层的 $\nabla_\theta \log \pi_\theta(A_t|S_t) Q^{\pi_\theta}(S_t,A_t)$ 
% % 只和 $S_t$ 和 $A_t$ 有关，所以期望的积分可以推前到 $S\times A$ 上的测度进行，即
% % \[
% %   \nabla_\theta J(\theta) = \sum_{t=0}^\infty \gamma^t\mathbb{E}_{s_t,a_t\sim \mathbb{P}_\theta}\left[
% %     \nabla_\theta \log \pi_\theta(a_t|s_t) Q^{\pi_\theta}(s_t,a_t)
% %   \right].
% % \]
% % 注意这里下标 $s_t,a_t\sim \mathbb{P}_\theta$ 是一种符号滥用，
% % 严格来说 $\mathbb{P}_\theta$ 应该推前到 $S\times A$ 上的测度 $(S_t,A_t)_*\mathbb{P}_\theta$。

% 为了降低策略梯度的方差，我们可以引入一个\emph{基线函数} $b:S\to \mathbb{R}$，
% 与 \eqref{eq:zero-mean} 类似，我们有
% \[
%   \mathbb{E}\left[
%     \strut\nabla_\theta \log \pi_\theta(A_t|S_t) b(S_t)\middle| S_t
%   \right] = b(S_t) \mathbb{E}\left[
%     \nabla_\theta \log \pi_\theta(A_t|S_t)\middle| S_t
%   \right] = 0.  
% \]
% 为了最小化估计方差，最好的选择是 $Q^{\pi_\theta}$ 的条件期望。定义\emph{状态价值函数}为
% \begin{align*}
%   V^{\pi_\theta}=\mathbb{E}[Q^{\pi_\theta}|S_t]=
%   \mathbb{E}\left[
%     \sum_{k=t}^\infty \gamma^{k-t} r_k
%     \middle| S_t
%   \right],
% \end{align*}
% 此时 $V^{\pi_\theta}$ 是 $S_t$ 的函数，表示在状态 $S_t$ 下未来累计奖励的条件期望。
% 差值 $A^{\pi_\theta}=Q^{\pi_\theta}-V^{\pi_\theta}$ 被称为\emph{优势函数}，表示在状态 $S_t$ 下执行动作 $A_t$ 相对于平均水平的优势。
% 由于
% \[
%   \mathbb{E}\left[
%     \nabla_\theta \log \pi_\theta(A_t|S_t) V^{\pi_\theta}(S_t)
%   \right]=\mathbb{E}\Bigl[
%   \mathbb{E}\bigl[
%     \nabla_\theta \log \pi_\theta(A_t|S_t) V^{\pi_\theta}(S_t)\big| S_t
%   \bigr]\Bigr]=0,
% \]
% 所以我们可以把策略梯度表达式改写成
% \begin{equation}\label{eq:policy-gradient-advantage}
%   \nabla_\theta J(\theta) = \sum_{t=0}^\infty \gamma^t\mathbb{E}_{\tau\sim \mathbb{P}_\theta}\left[
%     \nabla_\theta \log \pi_\theta(a_t|s_t) A^{\pi_\theta}(s_t,a_t)
%   \right].
% \end{equation}

% 记 $\nu_{t,\theta}$ 表示 $\mathbb{P}_\theta$ 在 $(S_t,A_t)$ 下的推前测度。
% $t=0$ 时，有
% \[
%   \nu_{0,\theta}(E\times F) = \mathbb{P}_\theta((S_0,A_0)\in E\times F) = \int_E \rho_0(\d s_0)\int_F \pi_\theta(\d a_0|s_0).
% \]
% $t=1$ 时，有
% \begin{align*}
%   \nu_{1,\theta}(E\times F) &= \int_{S\times A}\int_{E\times F} K_\theta(\d s_1,\d a_1|s_0,a_0)\rho_0(\d s_0)\pi_\theta(\d a_0|s_0)\\
%   &=\int_{S\times A}\int_{E\times F} \mathcal P(\d s_1|s_0,a_0)\pi_\theta(\d a_1|s_1)\rho_0(\d s_0)\pi_\theta(\d a_0|s_0)\\
%   &=\int_{S\times A}\nu_{0,\theta}(\d s_0\d a_0)\int_{E\times F} \mathcal P(\d s_1|s_0,a_0)\pi_\theta(\d a_1|s_1).
% \end{align*}
% 以此类推，对于一般的 $t$，我们有
% \begin{align*}
%   \nu_{t,\theta}(E\times F) &= \int_{S\times A}\nu_{t-1,\theta}(\d s_{t-1}\d a_{t-1})\int_{E\times F} \mathcal P(\d s_t|s_{t-1},a_{t-1})\pi_\theta(\d a_t|s_t)\\
%   &=\int_{S\times A}\nu_{t-1,\theta}(\d s_{t-1}\d a_{t-1})\int_E \mathcal P(\d s_t|s_{t-1},a_{t-1})
%   \pi_\theta(F|s_t)\\
%   &=\int_E \pi_\theta(F|s_t)\int_{S\times A}\nu_{t-1,\theta}(\d s_{t-1}\d a_{t-1})\mathcal P(\d s_t|s_{t-1},a_{t-1})\\
%   &=\int_E \pi_\theta(F|s_t)\mu_{t,\theta}(\d s_t)=
%   \int_{E\times F} \pi_\theta(\d a_t|s_t)\mu_{t,\theta}(\d s_t).
% \end{align*}
% 其中 
% \[
%   \mu_{t,\theta}(E) =\int_E \int_{S\times A}\nu_{t-1,\theta}(\d s_{t-1}\d a_{t-1})\mathcal P(\d s_t|s_{t-1},a_{t-1})
%   =\nu_{t,\theta}(E\times A)
% \]
% 是 $S$ 上的测度。

% 于是 \eqref{eq:policy-gradient-advantage} 中的积分可以推前到 $S\times A$ 上的测度 $\nu_{t,\theta}$，
% 得到
% \begin{align*}
%   &\mathbb{E}\left[
%     \nabla_\theta \log \pi_\theta(A_t|S_t) A^{\pi_\theta}(S_t,A_t)
%   \right] = \int_{S\times A} \nabla_\theta \log \pi_\theta(a|s) A^{\pi_\theta}(s,a) \nu_{t,\theta}(\d s \d a)\\
%   &=\int_S\mu_{t,\theta}(\d s)\int_A \nabla_\theta \log \pi_\theta(a|s) A^{\pi_\theta}(s,a) \pi_\theta(\d a|s),
% \end{align*}
% 交换积分和求和，得到
% \[
%   \nabla_\theta J(\theta)=\int_S \left(
%     \sum_{t=0}^\infty \gamma^t \mu_{t,\theta}(\d s)
%   \right)\int_A \nabla_\theta \log \pi_\theta(a|s) A^{\pi_\theta}(s,a) \pi_\theta(\d a|s).
% \]
% 定义\emph{归一化折现状态分布}为
% \[
%   \rho_{\pi_\theta}(E) = (1-\gamma)\sum_{t=0}^\infty \gamma^t \mu_{t,\theta}(E),
% \]
% 其是 $S$ 上的一个概率测度，表示在策略 $\pi_\theta$ 下系统访问状态集 $E$ 的归一化折现频率。
% 至此，我们得到了需要的\emph{最终策略梯度表达式}：
% \begin{equation}
%   \begin{aligned}
%     \nabla_\theta J(\theta) &= \frac{1}{1-\gamma}\int_S \rho_{\pi_\theta}(\d s)\int_A \nabla_\theta \log \pi_\theta(a|s) A^{\pi_\theta}(s,a) \pi_\theta(\d a|s)\\
%     &=\frac{1}{1-\gamma}\mathbb{E}_{s\sim\rho_{\pi_\theta},a\sim \pi_\theta}\left[
%       \nabla_\theta \log \pi_\theta(a|s) A^{\pi_\theta}(s,a)
%     \right].
%   \end{aligned}
% \end{equation}














\end{document}